% ============================================================
% 03_AUDITORIA.TEX
% Sistema Adaptativo de Classificação Multicritério
% de Fluidos de Corte
%
% FINALIDADE
%
% Apresentar:
% - por que a auditoria é necessária;
% - quais verificações serão realizadas;
% - como inconsistências serão tratadas;
% - como será construída a base canônica;
% - quais resultados da auditoria deverão ser apresentados.
%
% Este arquivo constitui apenas a BASE da apresentação.
%
% Nenhuma inconsistência específica deve ser apresentada como
% confirmada antes da execução da auditoria em R.
%
% ============================================================


% ============================================================
% FRAME 1 — POR QUE AUDITAR?
% ============================================================

\begin{frame}{Por que auditar os dados?}

    \begin{block}{Princípio}
        Um ranking multicritério só é confiável se os dados utilizados
        forem conhecidos, consistentes e rastreáveis.
    \end{block}

    \vspace{0.3cm}

    Antes de qualquer classificação, a auditoria deverá verificar:

    \begin{itemize}
        \item estrutura do arquivo;
        \item identificação das alternativas;
        \item tipos das variáveis;
        \item unidades;
        \item valores ausentes;
        \item duplicatas;
        \item fórmulas;
        \item consistência entre diferentes partes da base.
    \end{itemize}

    \vspace{0.2cm}

    \begin{alertblock}{Objetivo}
        Evitar que erros, divergências ou interpretações incorretas
        sejam propagados para CRITIC, TOPSIS e o ranking final.
    \end{alertblock}

\end{frame}


% ============================================================
% FRAME 2 — FLUXO DA AUDITORIA
% ============================================================

\begin{frame}{Fluxo da auditoria}

    \begin{block}{Processo previsto}
        \centering

        Inventário
        $\rightarrow$
        Validação estrutural
        $\rightarrow$
        Verificação de valores
        $\rightarrow$
        Comparação entre fontes
        $\rightarrow$
        Reconciliação
        $\rightarrow$
        Base canônica
    \end{block}

    \vspace{0.4cm}

    Cada problema identificado deverá possuir:

    \begin{itemize}
        \item localização;
        \item descrição;
        \item possível impacto;
        \item evidências;
        \item decisão;
        \item status.
    \end{itemize}

\end{frame}


% ============================================================
% FRAME 3 — VERIFICAÇÕES PRINCIPAIS
% ============================================================

\begin{frame}{Verificações principais}

    \begin{columns}[T]

        \begin{column}{0.48\textwidth}

            \begin{block}{Estrutura}
                \begin{itemize}
                    \item abas;
                    \item linhas e colunas;
                    \item identificadores;
                    \item tipos;
                    \item duplicatas;
                    \item campos ausentes.
                \end{itemize}
            \end{block}

            \begin{block}{Unidades}
                \begin{itemize}
                    \item unidade original;
                    \item consistência;
                    \item escala;
                    \item necessidade de conversão.
                \end{itemize}
            \end{block}

        \end{column}


        \begin{column}{0.48\textwidth}

            \begin{block}{Cálculos}
                \begin{itemize}
                    \item fórmulas;
                    \item somas;
                    \item médias;
                    \item totais;
                    \item índices derivados.
                \end{itemize}
            \end{block}

            \begin{block}{Consistência}
                \begin{itemize}
                    \item valores repetidos entre abas;
                    \item diferenças entre fontes;
                    \item erros de transcrição;
                    \item informações conflitantes.
                \end{itemize}
            \end{block}

        \end{column}

    \end{columns}

\end{frame}


% ============================================================
% FRAME 4 — VALORES AUSENTES E CASOS ESPECIAIS
% ============================================================

\begin{frame}{Valores ausentes e casos especiais}

    Nem toda ausência de informação possui o mesmo significado.

    \vspace{0.3cm}

    A auditoria deverá distinguir, quando possível:

    \begin{itemize}
        \item valor realmente ausente;
        \item medição pendente;
        \item informação não aplicável;
        \item erro de fórmula;
        \item conteúdo textual em campo numérico;
        \item informação de significado ainda desconhecido.
    \end{itemize}

    \vspace{0.3cm}

    \begin{alertblock}{Regra}
        Valores ausentes não serão convertidos automaticamente em zero.
    \end{alertblock}

\end{frame}


% ============================================================
% FRAME 5 — RECÁLCULO E RECONCILIAÇÃO
% ============================================================

\begin{frame}{Recálculo e reconciliação}

    \begin{block}{Recálculo independente}
        Valores derivados relevantes poderão ser recalculados em R para
        verificar a consistência da planilha original.
    \end{block}

    \vspace{0.3cm}

    Exemplos de verificações:

    \begin{itemize}
        \item somas;
        \item médias;
        \item percentuais;
        \item totais;
        \item índices;
        \item resultados derivados.
    \end{itemize}

    \vspace{0.3cm}

    Quando houver divergência:

    \begin{center}
        detectar
        $\rightarrow$
        investigar
        $\rightarrow$
        documentar
        $\rightarrow$
        decidir
        $\rightarrow$
        reconciliar
    \end{center}

\end{frame}


% ============================================================
% FRAME 6 — NÃO CORRIGIR SILENCIOSAMENTE
% ============================================================

\begin{frame}{Tratamento das inconsistências}

    \begin{alertblock}{Regra central}
        Nenhum valor será alterado apenas porque parece incorreto.
    \end{alertblock}

    \vspace{0.3cm}

    Uma inconsistência deverá ser resolvida com base em:

    \begin{itemize}
        \item origem do dado;
        \item significado técnico;
        \item unidade;
        \item cálculo independente;
        \item consistência com outras informações;
        \item documentação disponível.
    \end{itemize}

    \vspace{0.3cm}

    Caso não seja possível resolver:

    \begin{block}{Status}
        A pendência permanece explicitamente registrada e seu impacto nas
        etapas posteriores deverá ser avaliado.
    \end{block}

\end{frame}


% ============================================================
% FRAME 7 — BASE CANÔNICA
% ============================================================

\begin{frame}{Construção da base canônica}

    Após a auditoria e a reconciliação será gerada uma base canônica.

    \vspace{0.3cm}

    \begin{block}{Base canônica}
        Versão dos dados tratada, documentada e aprovada para utilização
        nas análises estatísticas posteriores.
    \end{block}

    \vspace{0.3cm}

    Ela deverá possuir:

    \begin{itemize}
        \item identificadores consistentes;
        \item variáveis documentadas;
        \item unidades definidas;
        \item transformações rastreáveis;
        \item inconsistências resolvidas ou registradas;
        \item versão identificável.
    \end{itemize}

\end{frame}


% ============================================================
% FRAME 8 — RASTREABILIDADE
% ============================================================

\begin{frame}{Rastreabilidade dos dados}

    Para cada transformação relevante, deverá ser possível reconstruir:

    \vspace{0.3cm}

    \begin{block}{Linhagem}
        \centering

        Valor analisado
        $\rightarrow$
        transformação
        $\rightarrow$
        decisão
        $\rightarrow$
        valor original
    \end{block}

    \vspace{0.4cm}

    Essa rastreabilidade permite:

    \begin{itemize}
        \item revisar decisões;
        \item reproduzir a análise;
        \item corrigir problemas na origem;
        \item evitar edições manuais não documentadas.
    \end{itemize}

\end{frame}


% ============================================================
% FRAME 9 — RESULTADOS DA AUDITORIA
% ============================================================

\begin{frame}{Resultados da auditoria}

    \begin{block}{A preencher após execução}

        \begin{itemize}

            \item Problemas identificados:
                  \texttt{TO\_BE\_FILLED}

            \item Problemas críticos:
                  \texttt{TO\_BE\_FILLED}

            \item Divergências reconciliadas:
                  \texttt{TO\_BE\_FILLED}

            \item Pendências:
                  \texttt{TO\_BE\_FILLED}

            \item Status da base canônica:
                  \texttt{TO\_BE\_FILLED}

        \end{itemize}

    \end{block}

    \vspace{0.3cm}

    \centering
    \textit{Essas informações deverão ser preenchidas por Pacheco
    somente após a auditoria em R.}

\end{frame}


% ============================================================
% FRAME 10 — ACHADOS MAIS IMPORTANTES
% ============================================================

\begin{frame}{Principais achados da auditoria}

    \begin{block}{Achado 1}
        \texttt{TO\_BE\_FILLED}
    \end{block}

    \begin{block}{Achado 2}
        \texttt{TO\_BE\_FILLED}
    \end{block}

    \begin{block}{Achado 3}
        \texttt{TO\_BE\_FILLED}
    \end{block}

    \vspace{0.2cm}

    \begin{alertblock}{Seleção para apresentação}
        Apresentar somente inconsistências que afetem de forma relevante
        os dados, a metodologia ou a interpretação final.
    \end{alertblock}

\end{frame}


% ============================================================
% FIGURAS FUTURAS
% ============================================================
%
% Depois da execução, esta seção poderá receber elementos como:
%
% - quadro resumo das inconsistências;
% - fluxograma da reconciliação;
% - comparação entre valor original e recalculado;
% - esquema da base canônica.
%
% Essas figuras deverão ser produzidas pelos integrantes e,
% quando científicas, preferencialmente geradas a partir do R.
%
% Não é necessário produzi-las agora.
%
% ============================================================


% ============================================================
% MATERIAL EXTERNO DA PROFESSORA
% ============================================================
%
% Caso a professora forneça explicações sobre:
%
% - origem de determinada variável;
% - significado de uma anotação;
% - unidade correta;
% - fórmula;
% - especificação;
%
% a informação deverá ser utilizada como evidência na
% reconciliação.
%
% Não substituir a informação oficial por inferência.
%
% ============================================================


% ============================================================
% ESTADO ATUAL
% ============================================================
%
% AUDIT_SECTION_STATUS = BASE_PREPARED
%
% AUDIT_STATUS = NOT_EXECUTED
% CANONICAL_DATA_STATUS = NOT_CREATED
% AUDIT_MAIN_FINDINGS = TO_BE_FILLED
%
% ============================================================